\documentclass[11pt]{article}
\usepackage[a4paper,margin=28mm,headheight=14pt]{geometry}
\usepackage[T1]{fontenc}
\usepackage{lmodern}
\usepackage{amsmath,amssymb,amsthm}
\usepackage{longtable,array}
\usepackage{fancyhdr}
\usepackage[hidelinks]{hyperref}

\hypersetup{pdftitle={Zhuk's Centers and Ternary Absorption},
  pdfsubject={A self-contained proof, at formalization granularity},pdfauthor={}}

\newtheoremstyle{mainplain}{9pt}{9pt}{\itshape}{}{\bfseries}{.}{0.55em}{}
\newtheoremstyle{mainroman}{9pt}{9pt}{\normalfont}{}{\bfseries}{.}{0.55em}{}
\theoremstyle{mainplain}
\newtheorem{theorem}{Theorem}[section]
\newtheorem{lemma}[theorem]{Lemma}
\newtheorem{proposition}[theorem]{Proposition}
\newtheorem{corollary}[theorem]{Corollary}
\theoremstyle{mainroman}
\newtheorem{longtheorem}[theorem]{Theorem}
\newtheorem{longlemma}[theorem]{Lemma}
\newtheorem{longproposition}[theorem]{Proposition}
\theoremstyle{definition}
\newtheorem{definition}[theorem]{Definition}
\newtheorem{remark}[theorem]{Remark}
\newtheorem{convention}[theorem]{Convention}

\setlength{\parindent}{1.25em}\setlength{\parskip}{0pt}\setlength{\emergencystretch}{2em}
\newcommand{\alphlabels}{\renewcommand{\labelenumi}{\textup{(\alph{enumi})}}%
  \renewcommand{\theenumi}{\alph{enumi}}}
\newcommand{\romanlabels}{\renewcommand{\labelenumi}{\textup{(\roman{enumi})}}%
  \renewcommand{\theenumi}{\roman{enumi}}}
\providecommand{\tightlist}{\setlength{\itemsep}{0pt}\setlength{\parskip}{0pt}}
\numberwithin{equation}{section}
\setcounter{tocdepth}{2}\setcounter{secnumdepth}{2}
\pagestyle{fancy}\fancyhf{}
\fancyhead[L]{\small Zhuk's centers}\fancyhead[R]{\small\nouppercase{\leftmark}}
\fancyfoot[C]{\thepage}\renewcommand{\headrulewidth}{0.3pt}

\DeclareMathOperator{\Sg}{Sg}
\DeclareMathOperator{\Clo}{Clo}
\newcommand{\bA}{\mathbf A}
\newcommand{\bB}{\mathbf B}
\newcommand{\bC}{\mathbf C}
\newcommand{\bD}{\mathbf D}
\newcommand{\bM}{\mathbf M}
\newcommand{\ab}{\lhd}
\newcommand{\abin}{\lhd_{\mathrm{bin}}}
\newcommand{\aZ}{\lhd_{Z}}
\newcommand{\Tm}{T_\sigma}

\title{Zhuk's Centers and Ternary Absorption}
\author{}\date{}

\begin{document}
\maketitle\vspace{-2.2em}
\begin{abstract}
Let $\bA$ and $\bB$ be finite idempotent algebras of a common signature and let
$R$ be a subdirect subuniverse of $\bA\times\bB$. The \emph{left center} of $R$
is the set $C$ of elements of $A$ related by $R$ to every element of $B$. We
give a complete and self-contained proof of Zhuk's theorem that if $\bB$ has a
Taylor term and has no proper binary absorbing subalgebra, then $C$ absorbs
$\bA$ --- and, through the Zhuk--Kozik essential doubling trick, that the
absorption is witnessed by a \emph{ternary} term. Everything is proved from the
definitions of terms, subuniverses, and generation: the relational description of
absorption by essential relations, which supplies the passage from unbounded
arity to ternary arity, is proved here rather than imported. The background
package of \hyperref[app:background]{Appendix~C} is finite combinatorics and
elementary set theory only. The document is written to serve as the
human-readable blueprint for a formal development.
\end{abstract}

\begin{theorem}[Main theorem]\label{thm:main}
Let $\sigma$ be a signature without nullary symbols and let $\bA,\bB$ be finite
idempotent $\sigma$-algebras with nonempty universes. Let $R\le\bA\times\bB$ be
subdirect and let
\[
  C=\{a\in A\;:\;\{a\}\times B\subseteq R\}
\]
be its left center. Assume
\begin{enumerate}\alphlabels\tightlist
\item some $t\in\Tm(k)$ is Taylor on $B$ \textup{(Definition~\ref{def:taylor})}, and
\item $\bB$ has no proper binary absorbing subalgebra
      \textup{(Definition~\ref{def:absorption})}.
\end{enumerate}
Then $C$ centrally absorbs $\bA$, and there is a \emph{ternary} term
$s\in\Tm(3)$ such that
\[
 s^{\bA}(C,A,C)\cup s^{\bA}(C,C,A)\cup s^{\bA}(A,C,C)\subseteq C .
\]
\end{theorem}

\tableofcontents\clearpage

\section*{Introduction}
\addcontentsline{toc}{section}{Introduction}

The algebraic approach to the constraint satisfaction problem studies a
finite relational template through the algebra of its polymorphisms. Zhuk's
proof of the dichotomy conjecture~\cite{Zhuk} isolated a mechanism that turns a
single binary relation between two algebras into an absorbing subalgebra of one
of them. Barto and Kozik's absorption theorem~\cite{BartoKozikCyclic} is the
older tool it replaces; Brady's notes~\cite{BradyNotes}, whose Section~3.10 this
document follows, use Zhuk's mechanism to reprove and strengthen it.

The mechanism is this. Given a subdirect $R\le\bA\times\bB$, call
$C=\{a: \{a\}\times B\subseteq R\}$ the \emph{left center}. If $C$ is neither
empty nor all of $A$, then $C$ is a nontrivial obstruction inside $\bA$, and one
expects to convert it into structure. Zhuk's observation is that the conversion
is forced: applying a Taylor term of $\bB$ to one element of $A$ and several
elements of $C$ strictly enlarges its $R$-neighborhood in $B$, unless that
neighborhood is a binary absorbing subalgebra of $\bB$. Iterating the
enlargement $|B|-1$ times drives the value into $C$. The witnessing term has
arity $k^{|B|-1}$; the second half of the argument, the Zhuk--Kozik doubling
trick, collapses that arity to $3$.

\medskip

\noindent\textbf{What is proved, and from what.} The document is
self-contained relative to \hyperref[app:background]{Appendix~C}, which asks
only for finite sets, induction, and elementary reasoning about functions and
products. In particular the following are proved rather than assumed:
\begin{itemize}\tightlist
\item that term operations preserve subuniverses, and the description of a
      generated subuniverse as the set of values of terms on a fixed list of
      generators (Section~\ref{sec:algebras});
\item the relational description of absorption --- $B$ absorbs $\bA$ with
      respect to an $m$-ary term precisely when no $m$-ary $B$-essential
      relation exists (Theorem~\ref{thm:relational-absorption}) --- which is the
      bridge between ``absorbs at some arity'' and ``absorbs at arity $3$'';
\item Zhuk's two theorems on left centers, the definition and basic theory of
      central absorption, the doubling trick, and the ternary conclusion
      (Section~\ref{sec:centers}).
\end{itemize}

\noindent\textbf{Route.} Part~I fixes the universal-algebraic vocabulary and
proves the handful of closure facts used everywhere: subuniverses are closed
under term operations, generated subuniverses are term images of a fixed
generator list, homomorphisms commute with generation, and the two relational
constructions (intersection with a box, and projection) preserve
subuniverses. It ends with absorption, binary absorption, and the lemma
--- Lemma~\ref{lem:taylor-binary} --- that converts a \emph{one-sided} closure
condition into binary absorption using a single Taylor identity.

Part~II develops essential relations. An $m$-ary $B$-essential relation is a
relation that meets every box with one coordinate free and all others confined
to $B$, but misses the all-$B$ box; it is exactly the obstruction to $m$-ary
absorption. The easy direction (absorption forbids essential relations) is a
one-line application of the witnessing term to the witnessing tuples. The hard
direction runs the free algebra $\Clo_m(\bA)$ against a grouped form of the
definition, and needs the regrouping Lemma~\ref{lem:essential-groups}.

Part~III is Section~3.10 of~\cite{BradyNotes}. Theorem~\ref{thm:center-step} is
the enlargement step, Theorem~\ref{thm:center-absorbs} its iteration,
Theorem~\ref{thm:center-central} the extra ``central'' property of the left
center, Lemma~\ref{lem:doubling} the doubling trick, and
Corollary~\ref{cor:central-ternary} the collapse to arity $3$.

\medskip

\noindent\textbf{Divergences from the source.} Three points where this document
is deliberately more explicit than~\cite{BradyNotes}. First, the empty set is a
subuniverse of every algebra in a signature without nullary symbols, and it
absorbs vacuously; the hypothesis ``$\bB$ has no proper binary absorbing
subalgebra'' must therefore be read as excluding only \emph{nonempty} proper
ones, which is what Convention~\ref{conv:subalgebra} records. Second,
Theorem~\ref{thm:center-central} is stated with the three generator blocks
$\{a\}\times C$, $C\times C$, $C\times\{a\}$; the middle block is genuinely
present in the source and is dropped only when passing to the weaker
Definition~\ref{def:central-absorption}. Third, the arities produced by
iterating the doubling trick from arity $3$ are $2+2^{k}$, and the argument in
Corollary~\ref{cor:central-ternary} is arranged here to need only that the
arity strictly increases.

\part{Algebras, terms, and absorption}
\label{part:algebras}

\section{Signatures, algebras, and subuniverses}\label{sec:algebras}

\begin{definition}[Signature and algebra]\label{def:algebra}
A \emph{signature} $\sigma$ is a set of operation symbols together with an
arity function $\mathrm{ar}:\sigma\to\{1,2,3,\dots\}$. Nullary symbols are not
permitted. A $\sigma$-\emph{algebra} $\bA$ consists of a set $A$, its
\emph{universe}, together with an operation
$f^{\bA}:A^{\mathrm{ar}(f)}\to A$ for each $f\in\sigma$. The algebra is
\emph{finite} if $A$ is finite. We write $\bA$ for the algebra and $A$ for its
universe, and never identify the two.

All algebras occurring below are $\sigma$-algebras for one fixed signature
$\sigma$, and all have nonempty universe.
\end{definition}

\begin{definition}[Idempotence]\label{def:idempotent}
$\bA$ is \emph{idempotent} if $f^{\bA}(a,a,\dots,a)=a$ for every $f\in\sigma$
and every $a\in A$.
\end{definition}

\begin{definition}[Subuniverse]\label{def:subuniverse}
A set $S\subseteq A$ is a \emph{subuniverse} of $\bA$, written $S\le\bA$, if
$f^{\bA}(s_1,\dots,s_{\mathrm{ar}(f)})\in S$ for every $f\in\sigma$ and all
$s_i\in S$. A subuniverse is \emph{proper} if it is not equal to $A$.
\end{definition}

Because $\sigma$ has no nullary symbols, $\varnothing\le\bA$ always. This is
convenient in Part~II, where projections of relations may be empty, but it makes
the phrase ``proper absorbing subalgebra'' ambiguous unless nonemptiness is
legislated. We legislate it once:

\begin{convention}[Subalgebra means nonempty]\label{conv:subalgebra}
A \emph{subalgebra} of $\bA$ is a nonempty subuniverse of $\bA$. Whenever a
hypothesis or conclusion says that an algebra \emph{has} or \emph{has no} proper
subalgebra with some property, ``subalgebra'' carries its nonemptiness. Objects
introduced as ``subuniverses'' may be empty unless stated otherwise.
\end{convention}

\begin{lemma}[Intersections and generation]\label{lem:generation}
Let $\bA$ be an algebra.
\begin{enumerate}\alphlabels\tightlist
\item The intersection of any nonempty family of subuniverses of $\bA$ is a
      subuniverse of $\bA$.
\item For $S\subseteq A$ there is a least subuniverse of $\bA$ containing $S$,
      namely the intersection of all subuniverses containing $S$. We write it
      $\Sg_{\bA}(S)$.
\item $S\subseteq T\le\bA$ implies $\Sg_{\bA}(S)\subseteq T$; and
      $S\subseteq S'$ implies $\Sg_{\bA}(S)\subseteq\Sg_{\bA}(S')$.
\end{enumerate}
\end{lemma}
\begin{proof}
For (a), if $s_1,\dots,s_n$ lie in every member of the family then so does
$f^{\bA}(s_1,\dots,s_n)$. For (b), the family of subuniverses containing $S$ is
nonempty because $A\le\bA$, so (a) applies; the intersection contains $S$ and is
contained in every subuniverse containing $S$. Part (c) is immediate from the
description in (b).
\end{proof}

\begin{lemma}[Singletons]\label{lem:singleton}
If $\bA$ is idempotent then $\{a\}\le\bA$ for every $a\in A$.
\end{lemma}
\begin{proof}
$f^{\bA}(a,\dots,a)=a$ by Definition~\ref{def:idempotent}.
\end{proof}

\begin{definition}[Products]\label{def:product}
Given algebras $\bA_1,\dots,\bA_m$, the product algebra
$\bA_1\times\cdots\times\bA_m$ has universe $A_1\times\cdots\times A_m$ and
operations applied coordinatewise. A \emph{relation} between
$\bA_1,\dots,\bA_m$ is a subuniverse $R\le\bA_1\times\cdots\times\bA_m$; its
\emph{arity} is $m$. We write $\bA^m$ for the $m$-fold product of $\bA$ with
itself, and $\bA^{X}$ for the product of copies of $\bA$ indexed by a set $X$,
whose universe is the set of functions $X\to A$.
\end{definition}

\begin{definition}[Terms]\label{def:term}
For $n\ge1$ let $\Tm(n)$ be the set of $n$-ary \emph{terms}, defined
inductively: the variables $x_1,\dots,x_n$ are $n$-ary terms, and if $f\in\sigma$
has arity $k$ and $t_1,\dots,t_k\in\Tm(n)$ then $f(t_1,\dots,t_k)\in\Tm(n)$.
For an algebra $\bA$, the \emph{term operation} $t^{\bA}:A^n\to A$ is defined by
$x_i^{\bA}(a_1,\dots,a_n)=a_i$ and
$f(t_1,\dots,t_k)^{\bA}(\vec a)=f^{\bA}\bigl(t_1^{\bA}(\vec a),\dots,t_k^{\bA}(\vec a)\bigr)$.
We write $\Clo_n(\bA)=\{t^{\bA}:t\in\Tm(n)\}\subseteq A^{A^n}$.

If $\pi:[n']\to[n]$ is any map and $t\in\Tm(n)$, the term
$t[\pi]\in\Tm(n')$ is obtained by replacing each variable $x_i$ by
$x_{\pi^{-1}}$; precisely, $t[\pi]$ is defined by structural recursion,
$x_i[\pi]=x_{j}$ for any chosen $j$ with $\pi(j)=i$. We use only the special
case where $\pi$ is a bijection, so that the choice is forced:
\emph{variable renaming along a permutation}. Then
$t[\pi]^{\bA}(a_1,\dots,a_n)=t^{\bA}(a_{\pi(1)},\dots,a_{\pi(n)})$.
\end{definition}

\begin{lemma}[Term operations are idempotent]\label{lem:term-idempotent}
If $\bA$ is idempotent then $t^{\bA}(a,a,\dots,a)=a$ for every $t\in\Tm(n)$ and
every $a\in A$.
\end{lemma}
\begin{proof}
Induction on $t$. For a variable the value is $a$. For $f(t_1,\dots,t_k)$ the
value is $f^{\bA}(a,\dots,a)=a$ by the inductive hypothesis and
Definition~\ref{def:idempotent}.
\end{proof}

\begin{lemma}[Preservation]\label{lem:preservation}
Let $S\le\bA$ and $t\in\Tm(n)$. Then $t^{\bA}(s_1,\dots,s_n)\in S$ for all
$s_i\in S$.
\end{lemma}
\begin{proof}
Induction on $t$. A variable returns one of the $s_i$. For $f(t_1,\dots,t_k)$,
each $t_j^{\bA}(\vec s)\in S$ by the inductive hypothesis, and $S$ is closed
under $f^{\bA}$ by Definition~\ref{def:subuniverse}.
\end{proof}

Applied to a product algebra, Lemma~\ref{lem:preservation} is the statement that
relations are \emph{preserved} by term operations, in the form used constantly
below: if $R\le\bA_1\times\cdots\times\bA_m$ and
$(a_{1j},\dots,a_{mj})\in R$ for $j=1,\dots,n$, then
\[
 \bigl(t^{\bA_1}(a_{11},\dots,a_{1n}),\;\dots,\;t^{\bA_m}(a_{m1},\dots,a_{mn})\bigr)\in R .
\]

\begin{lemma}[Generation by a fixed list]\label{lem:sg-terms}
Let $\bA$ be an algebra and let $S=\{s_1,\dots,s_N\}\subseteq A$ with $N\ge1$.
Then
\[
 \Sg_{\bA}(S)=\{\,t^{\bA}(s_1,\dots,s_N)\;:\;t\in\Tm(N)\,\}.
\]
\end{lemma}
\begin{proof}
Write $E$ for the right-hand side. Every element of $E$ lies in $\Sg_{\bA}(S)$
by Lemma~\ref{lem:preservation}, since $S\subseteq\Sg_{\bA}(S)\le\bA$.
Conversely $E$ contains each $s_i$, taking $t=x_i$; and $E\le\bA$, because for
$f\in\sigma$ of arity $k$ and $t_1,\dots,t_k\in\Tm(N)$,
\[
  f^{\bA}\bigl(t_1^{\bA}(\vec s),\dots,t_k^{\bA}(\vec s)\bigr)
  = \bigl(f(t_1,\dots,t_k)\bigr)^{\bA}(\vec s)\in E .
\]
Hence $\Sg_{\bA}(S)\subseteq E$ by Lemma~\ref{lem:generation}(c).
\end{proof}

\begin{remark}[Why a fixed list]\label{rem:fixed-list}
The usual formulation --- every element of $\Sg_{\bA}(S)$ is $t^{\bA}$ applied
to \emph{some} finite tuple of generators --- requires combining terms over
different generator lists, and hence a substitution operation on terms. Fixing
the list once, which is possible because all generating sets used below are
finite, avoids substitution entirely: only variable renaming along a permutation
is ever needed, and only in Theorem~\ref{thm:center-central}.
\end{remark}

\begin{definition}[Homomorphism]\label{def:hom}
A map $h:A\to B$ is a \emph{homomorphism} $\bA\to\bB$ if
$h(f^{\bA}(a_1,\dots,a_k))=f^{\bB}(h(a_1),\dots,h(a_k))$ for all $f\in\sigma$
and $a_i\in A$. Each coordinate projection
$\pi_i:\bA_1\times\cdots\times\bA_m\to\bA_i$, and more generally each
projection $\pi_J$ onto a subfamily of coordinates, is a homomorphism.
\end{definition}

\begin{lemma}[Images and generation]\label{lem:hom-image}
Let $h:\bA\to\bB$ be a homomorphism.
\begin{enumerate}\alphlabels\tightlist
\item If $S\le\bA$ then $h(S)\le\bB$.
\item $h(t^{\bA}(a_1,\dots,a_n))=t^{\bB}(h(a_1),\dots,h(a_n))$ for every
      $t\in\Tm(n)$.
\item $h(\Sg_{\bA}(S))=\Sg_{\bB}(h(S))$ for every finite nonempty
      $S\subseteq A$.
\end{enumerate}
\end{lemma}
\begin{proof}
(a) An element of $h(S)$ has the form $h(s)$; apply the homomorphism identity
and closure of $S$. (b) Induction on $t$, using the homomorphism identity at
each operation symbol. (c) Enumerate $S=\{s_1,\dots,s_N\}$; by
Lemma~\ref{lem:sg-terms} both sides equal
$\{t^{\bB}(h(s_1),\dots,h(s_N)):t\in\Tm(N)\}$, the left-hand side by (b).
\end{proof}

\begin{lemma}[Relational constructions]\label{lem:pp}
Let $\bA_1,\dots,\bA_m$ be algebras.
\begin{enumerate}\alphlabels\tightlist
\item If $R,R'\le\bA_1\times\cdots\times\bA_m$ then
      $R\cap R'\le\bA_1\times\cdots\times\bA_m$.
\item If $R\le\bA_1\times\cdots\times\bA_m$ and $S_i\le\bA_i$ for each $i$,
      then $R\cap(S_1\times\cdots\times S_m)\le\bA_1\times\cdots\times\bA_m$.
\item If $R\le\bA_1\times\cdots\times\bA_m$ and $J\subseteq[m]$ then
      $\pi_J(R)\le\prod_{i\in J}\bA_i$.
\item If $h:\bA\to\bB$ is a homomorphism and $S\le\bB$, then
      $h^{-1}(S)\le\bA$. In particular, for $J\subseteq[m]$ and
      $S\le\prod_{i\in J}\bA_i$, the cylinder $\pi_J^{-1}(S)$ is a subuniverse
      of $\bA_1\times\cdots\times\bA_m$.
\item If $\rho:[m']\to[m]$ is a bijection, the induced map
      $\vec a\mapsto(a_{\rho(1)},\dots,a_{\rho(m')})$ is an isomorphism
      $\bA_1\times\cdots\times\bA_m\to\bA_{\rho(1)}\times\cdots\times\bA_{\rho(m')}$,
      and carries subuniverses to subuniverses.
\end{enumerate}
\end{lemma}
\begin{proof}
(a) is Lemma~\ref{lem:generation}(a). (b) The box $S_1\times\cdots\times S_m$ is
a subuniverse because operations act coordinatewise, so (a) applies. (c) is
Lemma~\ref{lem:hom-image}(a) applied to the projection homomorphism. (d) If
$h(a_r)\in S$ for $r\in[k]$ then
$h(f^{\bA}(a_1,\dots,a_k))=f^{\bB}(h(a_1),\dots,h(a_k))\in S$. (e) The map is a
bijection commuting with coordinatewise operations, so
Lemma~\ref{lem:hom-image}(a) applies.
\end{proof}

\begin{remark}[The only relational constructions used]\label{rem:pp-usage}
Parts (a)--(e) of Lemma~\ref{lem:pp} are all that is needed to see that the
relations built in this document are subuniverses. They are used in exactly
three places: the neighborhoods $a+R$ of Definition~\ref{def:neighborhood}, the
grouped projections of Section~\ref{sec:essential}, and the doubled relation
$R'$ of Lemma~\ref{lem:doubling}. No general theory of primitive positive
definability is required.
\end{remark}

\section{Absorption}\label{sec:absorption}

\begin{definition}[Absorption]\label{def:absorption}
Let $\bA$ be an algebra and let $E\subseteq D$ be subuniverses of $\bA$. For
$t\in\Tm(m)$ we say that $t$ \emph{witnesses} $E\ab D$ if for every $i\in[m]$,
\[
 t^{\bA}(z_1,\dots,z_m)\in E
 \qquad\text{whenever } z_i\in D \text{ and } z_j\in E \text{ for all }j\ne i .
\]
We write $E\ab D$, and say $E$ \emph{absorbs} $D$, if some $t\in\Tm(m)$
witnesses it for some $m\ge1$; and $E\abin D$, and say $E$ \emph{binary
absorbs} $D$, if some $t\in\Tm(2)$ witnesses it, i.e.\ if
$t^{\bA}(E\times D)\cup t^{\bA}(D\times E)\subseteq E$. When $D=A$ we also
write $E\ab\bA$ and $E\abin\bA$.

The ambient algebra is suppressed in the notation. This is harmless: the
condition mentions $t^{\bA}$ only at arguments lying in $D$, so it depends on
$\bA$ only through the operations of $\bA$ restricted to $D$. No separate
notion of ``the subalgebra on $D$'' is needed.
\end{definition}

\begin{remark}[The quantifier form matters]\label{rem:absorption-quantifier}
Definition~\ref{def:absorption} constrains the tuple $(z_1,\dots,z_m)$ rather
than quantifying over a separate list $e_1,\dots,e_m\in E$ of which one entry is
then overwritten. The two readings differ exactly when $E=\varnothing$: the
second is vacuous for every $m\ge1$, while the first is vacuous only for
$m\ge2$ and \emph{fails} for $m=1$, where it asserts
$t^{\bA}(z_1)\in\varnothing$ for all $z_1\in D$. The first reading is the
correct one, in the sense that it is the one for which
Theorem~\ref{thm:relational-absorption} is true: for $S=\varnothing$ and $m=1$
the relation $R=A$ is $S$-essential of arity $1$, so the theorem must deny
$1$-ary absorption.
\end{remark}

\begin{remark}[Empty and full]\label{rem:absorption-degenerate}
$A\ab\bA$ holds, witnessed by $x_1$, and $\varnothing\ab\bA$ holds, witnessed
by any term of arity at least $2$. Convention~\ref{conv:subalgebra} is what
keeps the empty subuniverse from trivializing the hypothesis ``$\bB$ has no
proper binary absorbing subalgebra''.
\end{remark}

\begin{definition}[Taylor identities]\label{def:taylor}
Let $\bA$ be an algebra, $D\subseteq A$, $t\in\Tm(k)$ and $i\in[k]$. A
\emph{Taylor identity for $t$ at $i$ over $D$} is a pair of maps
$u,v:[k]\to\{1,2\}$ with $u(i)=1$ and $v(i)=2$ such that
\[
 t^{\bA}\bigl(w_{u(1)},\dots,w_{u(k)}\bigr)
 =t^{\bA}\bigl(w_{v(1)},\dots,w_{v(k)}\bigr)
 \qquad\text{for all }x,y\in D,
\]
where $w_1=x$ and $w_2=y$. We say $t$ is \emph{Taylor on $D$} if a Taylor
identity for $t$ at $i$ over $D$ exists for every $i\in[k]$.
\end{definition}

\begin{remark}[Relation to the usual definition]\label{rem:taylor-usual}
An algebra is usually called \emph{Taylor} if it has an idempotent term
satisfying, for each coordinate, an identity in two variables whose two sides
carry $x$ and $y$ respectively in that coordinate and arbitrary members of
$\{x,y\}$ elsewhere. Definition~\ref{def:taylor} is that condition, relativized
to a subset $D$ and stated for a term of the ambient signature; this is the
form in which it is used, since in Theorem~\ref{thm:center-step} the term is
applied to elements of $A$ while its identities are needed only on $B$.
\end{remark}

The next lemma is the engine that converts a one-sided closure condition into
two-sided binary absorption. It is the only place a Taylor identity is used.

\begin{lemma}[One-sided closure gives binary absorption]\label{lem:taylor-binary}
Let $\bA$ be an algebra, let $E\subseteq D\subseteq A$ with $D\le\bA$ and
$E\le\bA$, let $t\in\Tm(k)$ and $i\in[k]$, and suppose:
\begin{enumerate}\alphlabels\tightlist
\item there is a Taylor identity $(u,v)$ for $t$ at $i$ over $D$, and
\item $t^{\bA}(z_1,\dots,z_k)\in E$ whenever $z_i\in E$ and $z_j\in D$ for all
      $j\ne i$.
\end{enumerate}
Then $E\abin D$ in the sense of Definition~\ref{def:absorption}; explicitly, the
binary term $f:=t\bigl(x_{u(1)},\dots,x_{u(k)}\bigr)\in\Tm(2)$ satisfies
$f^{\bA}(E\times D)\cup f^{\bA}(D\times E)\subseteq E$.
\end{lemma}
\begin{proof}
Let $x,y\in D$ and put $w_1=x$, $w_2=y$. By construction
$f^{\bA}(x,y)=t^{\bA}(w_{u(1)},\dots,w_{u(k)})$, and by (a) this also equals
$t^{\bA}(w_{v(1)},\dots,w_{v(k)})$.

Take $e\in E$ and $d\in D$. For $f^{\bA}(e,d)$ set $x=e$, $y=d$ and use the
$u$-form: its $i$th argument is $w_{u(i)}=w_1=e\in E$, and every argument is one
of $e,d$, hence lies in $D$ since $E\subseteq D$. By (b),
$f^{\bA}(e,d)\in E$. For $f^{\bA}(d,e)$ set $x=d$, $y=e$ and use the $v$-form:
its $i$th argument is $w_{v(i)}=w_2=e\in E$, and again every argument lies in
$D$. By (b), $f^{\bA}(d,e)\in E$.
\end{proof}

\begin{definition}[Star powers]\label{def:star}
Let $t\in\Tm(k)$ with $k\ge1$. Define $t^{*0}:=x_1\in\Tm(1)$ and, for
$\ell\ge0$,
\[
 t^{*(\ell+1)}:=t\Bigl(
   t^{*\ell}\bigl(x_1,\dots,x_{k^{\ell}}\bigr),\;
   t^{*\ell}\bigl(x_{k^{\ell}+1},\dots,x_{2k^{\ell}}\bigr),\;\dots,\;
   t^{*\ell}\bigl(x_{(k-1)k^{\ell}+1},\dots,x_{k^{\ell+1}}\bigr)\Bigr)
 \;\in\;\Tm\bigl(k^{\ell+1}\bigr).
\]
Writing $p\in[k^{\ell+1}]$ uniquely as $p=(j-1)k^{\ell}+q$ with $j\in[k]$ and
$q\in[k^{\ell}]$, we call $j$ the \emph{block} of $p$ and $q$ its
\emph{position in the block}, so that
\[
 \bigl(t^{*(\ell+1)}\bigr)^{\bA}(a_1,\dots,a_{k^{\ell+1}})
 = t^{\bA}\bigl(d_1,\dots,d_k\bigr),\qquad
 d_j=\bigl(t^{*\ell}\bigr)^{\bA}\bigl(a_{(j-1)k^{\ell}+1},\dots,a_{jk^{\ell}}\bigr).
\]
\end{definition}

\part{Essential relations and the arity of absorption}
\label{part:essential}

\section{Essential relations}\label{sec:essential}

\begin{definition}[Essential relation]\label{def:essential}
Let $\bA_1,\dots,\bA_m$ be algebras with $m\ge1$ and let $S_i\le\bA_i$. A
relation $R\le\bA_1\times\cdots\times\bA_m$ is
$(S_1,\dots,S_m)$-\emph{essential} if
\begin{align}
 &R\cap(S_1\times\cdots\times S_{i-1}\times A_i\times S_{i+1}\times\cdots\times S_m)\ne\varnothing
 \quad\text{for every }i\in[m],\label{eq:E1}\\
 &R\cap(S_1\times\cdots\times S_m)=\varnothing.\label{eq:E2}
\end{align}
When all $\bA_i$ equal $\bA$ and all $S_i$ equal $S$, we say $R$ is
$S$-\emph{essential of arity} $m$.
\end{definition}

\begin{lemma}[Essentiality forces nonemptiness]\label{lem:essential-nonempty}
If $R\le\bA_1\times\cdots\times\bA_m$ is $(S_1,\dots,S_m)$-essential and
$m\ge2$, then $S_i\ne\varnothing$ for every $i\in[m]$.
\end{lemma}
\begin{proof}
Fix $i$ and choose $i'\ne i$ in $[m]$. The witness supplied by
\eqref{eq:E1} for the index $i'$ is a tuple whose $i$th coordinate lies in
$S_i$.
\end{proof}

\begin{proposition}[Passing to smaller arity]\label{prop:essential-down}
Let $R\le\bA_1\times\cdots\times\bA_m$ be $(S_1,\dots,S_m)$-essential with
$m\ge2$. Then
\[
 R^{-}:=\pi_{[m-1]}\bigl(R\cap(A_1\times\cdots\times A_{m-1}\times S_m)\bigr)
\]
is an $(S_1,\dots,S_{m-1})$-essential relation
$R^{-}\le\bA_1\times\cdots\times\bA_{m-1}$. Consequently, if an
$S$-essential relation of arity $m$ exists, then $S$-essential relations of
every arity $m'$ with $1\le m'\le m$ exist.
\end{proposition}
\begin{proof}
$R^{-}$ is a subuniverse by Lemma~\ref{lem:pp}(b),(c).

\eqref{eq:E1}: fix $i\le m-1$ and let $\vec x$ be the witness for index $i$
supplied by \eqref{eq:E1} for $R$. Its last coordinate lies in $S_m$, so
$\vec x\in R\cap(A_1\times\cdots\times A_{m-1}\times S_m)$, and
$\pi_{[m-1]}(\vec x)$ lies in the required box.

\eqref{eq:E2}: suppose $\vec y\in R^{-}\cap(S_1\times\cdots\times S_{m-1})$.
Then $\vec y=\pi_{[m-1]}(\vec x)$ for some $\vec x\in R$ with $x_m\in S_m$,
whence $\vec x\in R\cap(S_1\times\cdots\times S_m)$, contradicting
\eqref{eq:E2} for $R$.

The final sentence follows by iterating, each step reducing the arity by one.
\end{proof}

\begin{proposition}[Absorption forbids essential relations]\label{prop:absorption-blocks}
Let $S\le\bA$ and suppose $t\in\Tm(m)$ witnesses $S\ab\bA$. Then there is no
$S$-essential relation of arity $m$.
\end{proposition}
\begin{proof}
Suppose $R\le\bA^m$ is $S$-essential. For each $i\in[m]$ let $\vec r^{\,(i)}\in R$
be the witness supplied by \eqref{eq:E1}, so $r^{(i)}_i\in A$ and
$r^{(i)}_j\in S$ for $j\ne i$. Apply $t^{\bA^m}$ to the $m$ tuples
$\vec r^{\,(1)},\dots,\vec r^{\,(m)}$; by Lemma~\ref{lem:preservation} the result
$\vec z$ lies in $R$, and its $j$th coordinate is
\[
  z_j=t^{\bA}\bigl(r^{(1)}_j,\dots,r^{(m)}_j\bigr)
     =t^{\bA}\bigl(\underbrace{r^{(1)}_j}_{\in S},\dots,
        \underbrace{r^{(j)}_j}_{\in A},\dots,\underbrace{r^{(m)}_j}_{\in S}\bigr)\in S
\]
by Definition~\ref{def:absorption}, the free argument sitting in position $j$.
Hence $\vec z\in R\cap S^m$, contradicting \eqref{eq:E2}.
\end{proof}

\begin{corollary}[Absorption bounds essential arity]\label{cor:absorption-bound}
If $S\ab\bA$ is witnessed by a term of arity $m$, then every $S$-essential
relation has arity at most $m-1$. In particular the set of arities of
$S$-essential relations is bounded.
\end{corollary}
\begin{proof}
An $S$-essential relation of arity $m'\ge m$ would yield one of arity exactly
$m$ by Proposition~\ref{prop:essential-down}, contradicting
Proposition~\ref{prop:absorption-blocks}.
\end{proof}

The converse of Proposition~\ref{prop:absorption-blocks} is the substantial
result of this part. It goes through a grouped form of essentiality.

\begin{lemma}[Regrouping]\label{lem:essential-groups}
Let $S\le\bA$, let $m\ge1$ and $n_1,\dots,n_m\ge1$, and regard
\[
 R\;\le\;\bA^{n_1}\times\cdots\times\bA^{n_m}
\]
as a relation of arity $n=n_1+\cdots+n_m$ whose coordinates are partitioned into
$m$ consecutive blocks of sizes $n_1,\dots,n_m$. If $R$ is
$(S^{n_1},\dots,S^{n_m})$-essential --- that is, essential for the block
partition, with $S^{n_j}$ the box on block $j$ --- then there exists an
$S$-essential relation of arity $m$.
\end{lemma}
\begin{proof}
By strong induction on $n$, with $m$ fixed.

If $n_1=\cdots=n_m=1$ then $n=m$ and $R$ is itself an $S$-essential relation of
arity $m$.

Otherwise some $n_j>1$. Permuting the blocks permutes the coordinates by a
block permutation, which by Lemma~\ref{lem:pp}(e) carries $R$ to a subuniverse
of the correspondingly permuted product, and permutes the two essentiality
conditions among themselves; so we may assume $n_m>1$. Write
$n_{<m}=n_1+\cdots+n_{m-1}$, so $n=n_{<m}+n_m$ and the coordinates of the last
block are $n_{<m}+1,\dots,n$.

Consider
\[
 R_1:=\pi_{[n-1]}\bigl(R\cap(A^{n-1}\times S)\bigr)
 \;\le\;\bA^{n_1}\times\cdots\times\bA^{n_{m-1}}\times\bA^{n_m-1},
\]
a relation of arity $n-1$ with block sizes $n_1,\dots,n_{m-1},n_m-1$, all
$\ge1$. It is a subuniverse by Lemma~\ref{lem:pp}(b),(c). We check the two
essentiality conditions for these blocks.

\emph{Condition \eqref{eq:E2}.} If $\vec y\in R_1$ has all coordinates in $S$
then $\vec y=\pi_{[n-1]}(\vec x)$ with $\vec x\in R$ and $x_n\in S$, so
$\vec x\in R\cap S^n$, contradicting \eqref{eq:E2} for $R$.

\emph{Condition \eqref{eq:E1} for blocks $j\le m-1$.} The witness $\vec x\in R$
for block $j$ has all coordinates outside block $j$ in $S$; in particular
$x_n\in S$, so $\vec x\in R\cap(A^{n-1}\times S)$, and $\pi_{[n-1]}(\vec x)$ is
a witness for block $j$ in $R_1$.

\emph{Condition \eqref{eq:E1} for the last block.} This asks for
$\vec x\in R$ with $x_1,\dots,x_{n_{<m}}\in S$ and $x_n\in S$, the remaining
coordinates $x_{n_{<m}+1},\dots,x_{n-1}$ unconstrained; equivalently, it asks
that
\[
 \pi_{[n_{<m}]\cup\{n\}}(R)\cap S^{\,n_{<m}+1}\ne\varnothing. \tag{$\dagger$}
\]

If $(\dagger)$ holds then $R_1$ is essential for its block partition, of arity
$n-1<n$, and the inductive hypothesis applies.

If $(\dagger)$ fails, put
\[
 R_2:=\pi_{[n_{<m}]\cup\{n\}}(R)\;\le\;\bA^{n_1}\times\cdots\times\bA^{n_{m-1}}\times\bA,
\]
a relation of arity $n_{<m}+1$ with block sizes $n_1,\dots,n_{m-1},1$, all
$\ge1$; it is a subuniverse by Lemma~\ref{lem:pp}(c). Condition \eqref{eq:E2}
for $R_2$ is precisely the failure of $(\dagger)$. For \eqref{eq:E1} with
$j\le m-1$, project the witness of $R$ for block $j$: its coordinates outside
block $j$ lie in $S$, in particular $x_n\in S$. For \eqref{eq:E1} on the last
block, project the witness of $R$ for block $m$: its first $n_{<m}$ coordinates
lie in $S$, and its $n$th coordinate is unconstrained. Since $n_m>1$ we have
$n_{<m}+1<n$, so the inductive hypothesis applies to $R_2$.
\end{proof}

\begin{lemma}[The free algebra on $m$ generators]\label{lem:clo-subuniverse}
For every $m\ge1$, $\Clo_m(\bA)$ is a subuniverse of $\bA^{A^m}$ containing the
$m$ coordinate projections $p_1,\dots,p_m:A^m\to A$.
\end{lemma}
\begin{proof}
$p_i=x_i^{\bA}\in\Clo_m(\bA)$. For $f\in\sigma$ of arity $k$ and
$t_1,\dots,t_k\in\Tm(m)$, the pointwise application
$f^{\bA^{A^m}}(t_1^{\bA},\dots,t_k^{\bA})$ sends $\vec a\mapsto
f^{\bA}(t_1^{\bA}(\vec a),\dots,t_k^{\bA}(\vec a))
=\bigl(f(t_1,\dots,t_k)\bigr)^{\bA}(\vec a)$, so it lies in $\Clo_m(\bA)$.
\end{proof}

\begin{theorem}[Relational description of absorption]\label{thm:relational-absorption}
Let $\bA$ be a finite algebra, $S\le\bA$, and $m\ge1$. Then $S\ab\bA$ is
witnessed by some term of arity $m$ if and only if there is no $S$-essential
relation of arity $m$.
\end{theorem}
\begin{proof}
The forward implication is Proposition~\ref{prop:absorption-blocks}.

For the converse, assume there is no $S$-essential relation of arity $m$.

If $S=A$, every $t\in\Tm(m)$ witnesses $S\ab\bA$, since all its values lie in
$A=S$; and $x_1\in\Tm(m)$, so such a $t$ exists.

If $S=\varnothing$ and $m\ge2$, the witnessing condition of
Definition~\ref{def:absorption} is vacuous, because it requires $z_j\in S$ for
the $m-1\ge1$ indices $j\ne i$. If $S=\varnothing$ and $m=1$, the hypothesis is
unsatisfiable: $R=A$ is an $S$-essential relation of arity $1$, being nonempty
and disjoint from $S$. So in either case there is nothing to prove.

Assume from now on that $\varnothing\ne S\subsetneq A$.

For $i\in[m]$ set
\[
 X_i:=\{\vec x\in A^m\;:\;x_i\in A\setminus S\text{ and }x_j\in S\text{ for }j\ne i\}.
\]
Each $X_i$ is nonempty, since $S$ and $A\setminus S$ are nonempty. The $X_i$ are
pairwise disjoint, because a tuple in $X_i\cap X_{i'}$ with $i\ne i'$ would have
$x_i\in S$ and $x_i\notin S$. Put $X:=X_1\cup\cdots\cup X_m$ and
\[
 R:=\pi_{X}\bigl(\Clo_m(\bA)\bigr)\;\le\;\bA^{X_1}\times\cdots\times\bA^{X_m},
\]
using Lemma~\ref{lem:clo-subuniverse} and Lemma~\ref{lem:pp}(c). Concretely, $R$
is the set of restrictions $t^{\bA}\!\restriction_X$ for $t\in\Tm(m)$, and we
regard it as a relation of arity $|X|$ with the $m$ blocks $X_1,\dots,X_m$, of
sizes $n_i=|X_i|\ge1$.

\emph{Condition \eqref{eq:E1}.} Fix $i\in[m]$ and take $t=x_i$, so
$t^{\bA}=p_i$. For $\vec x\in X_j$ with $j\ne i$ we have $x_i\in S$ by the
definition of $X_j$, so $p_i(\vec x)\in S$. Hence
$\pi_X(p_i)$ lies in $S^{X_1}\times\cdots\times A^{X_i}\times\cdots\times S^{X_m}$,
witnessing \eqref{eq:E1} for block $i$.

Since no $S$-essential relation of arity $m$ exists, Lemma~\ref{lem:essential-groups}
shows that $R$ is not $(S^{X_1},\dots,S^{X_m})$-essential. Condition
\eqref{eq:E1} holds, so \eqref{eq:E2} must fail:
\[
 R\cap\bigl(S^{X_1}\times\cdots\times S^{X_m}\bigr)\ne\varnothing .
\]
Choose $t\in\Tm(m)$ whose restriction to $X$ witnesses this, so
$t^{\bA}(\vec x)\in S$ for every $\vec x\in X$.

Finally, $t$ witnesses $S\ab\bA$. Let $i\in[m]$, $a\in A$, $s_1,\dots,s_m\in S$
and put $\vec x=(s_1,\dots,s_{i-1},a,s_{i+1},\dots,s_m)$. If $a\notin S$ then
$\vec x\in X_i$ and $t^{\bA}(\vec x)\in S$ as just shown. If $a\in S$ then
$\vec x\in S^m$ and $t^{\bA}(\vec x)\in S$ by Lemma~\ref{lem:preservation}.
\end{proof}

\begin{corollary}[Unbounded arity defeats absorption]\label{cor:unbounded-no-absorption}
Let $\bA$ be a finite algebra and $S\le\bA$ with $S\ab\bA$. Then there is a
largest arity of an $S$-essential relation, or no $S$-essential relation at all.
\end{corollary}
\begin{proof}
Immediate from Corollary~\ref{cor:absorption-bound}: the arities form a set of
positive integers bounded above, hence finite, hence with a largest element if
nonempty.
\end{proof}

\part{Zhuk's centers}
\label{part:centers}

\section{Left centers and their neighborhoods}\label{sec:centers}

Throughout this part $\bA$ and $\bB$ are finite idempotent algebras with
nonempty universes, and $R\le\bA\times\bB$.

\begin{definition}[Subdirect]\label{def:subdirect}
$R\le\bA\times\bB$ is \emph{subdirect} if $\pi_1(R)=A$ and $\pi_2(R)=B$.
\end{definition}

\begin{definition}[Neighborhood and left center]\label{def:neighborhood}
For $a\in A$ put
\[
 a+R:=\{b\in B\;:\;(a,b)\in R\}\;=\;\pi_2\bigl(R\cap(\{a\}\times B)\bigr),
\]
and define the \emph{left center} of $R$ to be
\[
 C:=\{a\in A\;:\;a+R=B\}=\{a\in A: \{a\}\times B\subseteq R\}.
\]
The \emph{right center} of $R$ is the left center of
$R^{-}=\{(b,a):(a,b)\in R\}$, a subuniverse of $\bB\times\bA$.
\end{definition}

\begin{lemma}[Basic properties]\label{lem:center-basic}
Let $R\le\bA\times\bB$ be subdirect with left center $C$.
\begin{enumerate}\alphlabels\tightlist
\item $a+R\le\bB$ for every $a\in A$, and $a+R\ne\varnothing$.
\item $|a+R|\le|B|$, with equality if and only if $a\in C$.
\item $C\le\bA$.
\item For every $t\in\Tm(k)$ and $a_1,\dots,a_k\in A$,
      \[
        t^{\bA}(a_1,\dots,a_k)+R\;\supseteq\;
        \bigl\{\,t^{\bB}(b_1,\dots,b_k)\;:\;b_j\in a_j+R\,\bigr\}.
      \]
\end{enumerate}
\end{lemma}
\begin{proof}
(a) $\{a\}\le\bA$ by Lemma~\ref{lem:singleton}, so $a+R\le\bB$ by
Lemma~\ref{lem:pp}(b),(c). It is nonempty because $\pi_1(R)=A$.

(b) Immediate from $a+R\subseteq B$ and the definition of $C$.

(c) Let $f\in\sigma$ have arity $k$, let $a_1,\dots,a_k\in C$ and let $b\in B$.
Then $(a_j,b)\in R$ for every $j$, so
$(f^{\bA}(a_1,\dots,a_k),f^{\bB}(b,\dots,b))\in R$; and
$f^{\bB}(b,\dots,b)=b$ because $\bB$ is idempotent. As $b\in B$ was arbitrary,
$f^{\bA}(a_1,\dots,a_k)\in C$.

(d) If $b_j\in a_j+R$ then $(a_j,b_j)\in R$ for each $j$, so
$(t^{\bA}(\vec a),t^{\bB}(\vec b))\in R$ by Lemma~\ref{lem:preservation}.
\end{proof}

\begin{remark}[Degenerate centers]\label{rem:degenerate-center}
If $C=\varnothing$ then every assertion below about $C$ holds vacuously
(Remark~\ref{rem:absorption-degenerate}); if $C=A$ then $R=A\times B$ and the
assertions are trivial. The content is in the case
$\varnothing\ne C\subsetneq A$, but no proof needs to exclude the degenerate
cases, and none does.
\end{remark}

\section{The enlargement step}\label{sec:step}

\begin{theorem}[One step towards the center]\label{thm:center-step}
Let $\bA,\bB$ be finite idempotent algebras, let $R\le\bA\times\bB$ be
subdirect with left center $C$, and let $t\in\Tm(k)$ be Taylor on $B$. Assume
$\bB$ has no proper binary absorbing subalgebra. Then for every $a\in A$, every
$i\in[k]$ and all $c_1,\dots,c_k\in C$,
\[
 \bigl|\,t^{\bA}(c_1,\dots,c_{i-1},a,c_{i+1},\dots,c_k)+R\,\bigr|
 \;\ge\;\min\bigl(|B|,\;|a+R|+1\bigr).
\]
\end{theorem}
\begin{proof}
Write $a^{+}:=t^{\bA}(c_1,\dots,a,\dots,c_k)$ with $a$ in position $i$, and set
\[
 E:=\bigl\{\,t^{\bB}(z_1,\dots,z_k)\;:\;z_i\in a+R,\;z_j\in B\ (j\ne i)\,\bigr\}.
\]

\emph{Step 1: $a^{+}+R\supseteq E$.} Each $c_j\in C$ has $c_j+R=B$, so
Lemma~\ref{lem:center-basic}(d) gives exactly this containment.

\emph{Step 2: $E\supseteq a+R$.} For $b\in a+R$ take $z_1=\cdots=z_k=b$; this is
admissible because $a+R\subseteq B$, and $t^{\bB}(b,\dots,b)=b$ by
Lemma~\ref{lem:term-idempotent}.

\emph{Step 3: the dichotomy.} If $E\ne a+R$ then $|E|\ge|a+R|+1$ by Step 2, and
$|a^{+}+R|\ge|E|\ge|a+R|+1$ by Step 1, which is the claim.

So suppose $E=a+R$. Then $t^{\bB}(z_1,\dots,z_k)\in a+R$ whenever
$z_i\in a+R$ and $z_j\in B$ for $j\ne i$. Apply Lemma~\ref{lem:taylor-binary}
inside $\bB$, with $D=B$, with $E_{\text{lem}}=a+R$ --- a subuniverse of $\bB$
by Lemma~\ref{lem:center-basic}(a) --- and with the Taylor identity for $t$ at
$i$ over $B$ supplied by Definition~\ref{def:taylor}. It yields
$a+R\abin\bB$. Now $a+R$ is nonempty by Lemma~\ref{lem:center-basic}(a), so by
the hypothesis on $\bB$ and Convention~\ref{conv:subalgebra} it is not a proper
subuniverse: $a+R=B$. Hence $a\in C$, so $a^{+}\in C$ by
Lemma~\ref{lem:center-basic}(c), so $|a^{+}+R|=|B|\ge\min(|B|,|a+R|+1)$.
\end{proof}

\begin{theorem}[Zhuk: the left center absorbs]\label{thm:center-absorbs}
Let $\bA,\bB$ be finite idempotent algebras, let $R\le\bA\times\bB$ be
subdirect with left center $C$, and let $t\in\Tm(k)$ be Taylor on $B$. If
$\bB$ has no proper binary absorbing subalgebra, then the term
$t^{*(|B|-1)}$ of Definition~\ref{def:star} witnesses $C\ab\bA$. In particular
$C\ab\bA$.
\end{theorem}
\begin{proof}
We prove by induction on $\ell\ge0$ the following statement: for every
$a\in A$, every $p\in[k^{\ell}]$ and all $c_1,\dots,c_{k^{\ell}}\in C$,
\[
 \bigl|\,\bigl(t^{*\ell}\bigr)^{\bA}(c_1,\dots,c_{p-1},a,c_{p+1},\dots,c_{k^{\ell}})+R\,\bigr|
 \;\ge\;\min\bigl(|B|,\;|a+R|+\ell\bigr).
 \tag{$\ast_\ell$}
\]

For $\ell=0$ the term is $x_1$ and the value is $a$, so ($\ast_0$) reads
$|a+R|\ge\min(|B|,|a+R|)$, which holds by
Lemma~\ref{lem:center-basic}(b).

Assume ($\ast_\ell$). Let $a\in A$, $p\in[k^{\ell+1}]$ and
$c_1,\dots,c_{k^{\ell+1}}\in C$, and write $p=(j-1)k^{\ell}+q$ with $j\in[k]$,
$q\in[k^{\ell}]$ as in Definition~\ref{def:star}. Then
\[
 \bigl(t^{*(\ell+1)}\bigr)^{\bA}(c_1,\dots,a,\dots,c_{k^{\ell+1}})
 = t^{\bA}(d_1,\dots,d_k),
\]
where $d_j=\bigl(t^{*\ell}\bigr)^{\bA}$ applied to the $j$th block, which
consists of elements of $C$ except for $a$ in position $q$, and where for
$j'\ne j$ the value $d_{j'}=\bigl(t^{*\ell}\bigr)^{\bA}$ is applied to a block
of elements of $C$ only, so $d_{j'}\in C$ by
Lemma~\ref{lem:center-basic}(c) and Lemma~\ref{lem:preservation}.

By ($\ast_\ell$), $|d_j+R|\ge\min(|B|,|a+R|+\ell)$. By
Theorem~\ref{thm:center-step} applied to $d_j$ in position $j$ with the
$C$-elements $d_{j'}$,
\[
 \bigl|t^{\bA}(d_1,\dots,d_k)+R\bigr|\ \ge\ \min\bigl(|B|,\;|d_j+R|+1\bigr)
 \ \ge\ \min\bigl(|B|,\;|a+R|+\ell+1\bigr),
\]
the last step because $\min(|B|,\min(|B|,x)+1)=\min(|B|,x+1)$ for
$x\ge0$. This is ($\ast_{\ell+1}$).

Now take $\ell=|B|-1$ and let $a\in A$, $p\in[k^{\ell}]$,
$c_1,\dots,c_{k^{\ell}}\in C$. Since $|a+R|\ge1$ by
Lemma~\ref{lem:center-basic}(a), ($\ast_\ell$) gives
\[
 \bigl|\bigl(t^{*\ell}\bigr)^{\bA}(c_1,\dots,a,\dots,c_{k^{\ell}})+R\bigr|
 \ \ge\ \min\bigl(|B|,\,1+|B|-1\bigr)=|B|,
\]
so the value lies in $C$ by Lemma~\ref{lem:center-basic}(b). As $p$ was
arbitrary, $t^{*\ell}$ witnesses $C\ab\bA$ in the sense of
Definition~\ref{def:absorption}.
\end{proof}

\begin{remark}[The degenerate arity]\label{rem:star-degenerate}
If $|B|=1$ the exponent is $\ell=0$ and the witnessing term is $x_1$; this is
correct, since $R$ subdirect and $|B|=1$ force $R=A\times B$ and $C=A$.
\end{remark}

\section{The center is central}\label{sec:central}

\begin{theorem}[Zhuk: centrality of the left center]\label{thm:center-central}
Let $\bA,\bB$ be finite idempotent algebras and let $R\le\bA\times\bB$ be
subdirect with left center $C$. Assume $\bB$ has no proper binary absorbing
subalgebra. Then for every $a\in A\setminus C$,
\[
 (a,a)\;\notin\;\Sg_{\bA^2}\Bigl(\bigl(\{a\}\times C\bigr)\cup\bigl(C\times C\bigr)\cup\bigl(C\times\{a\}\bigr)\Bigr).
\]
\end{theorem}
\begin{proof}
Write $G:=(\{a\}\times C)\cup(C\times C)\cup(C\times\{a\})$ and suppose
$(a,a)\in\Sg_{\bA^2}(G)$. If $C=\varnothing$ then $G=\varnothing$ and
$\Sg_{\bA^2}(\varnothing)=\varnothing$, since $\sigma$ has no nullary symbols;
so $C\ne\varnothing$ and $G\ne\varnothing$.

\emph{Normal form for the generators.} The three sets $\{a\}\times C$,
$C\times C$ and $C\times\{a\}$ are pairwise disjoint: a pair lying in two of
them would force $a\in C$. Enumerate $G$ so that the elements of
$\{a\}\times C$ come first, then those of $C\times C$, then those of
$C\times\{a\}$. By Lemma~\ref{lem:sg-terms} there are $N\ge1$, an enumeration
$g_1,\dots,g_N$ of $G$ of this shape and a term $s\in\Tm(N)$ with
$(a,a)=s^{\bA^2}(g_1,\dots,g_N)$. Let $n:=N$, $t:=s$ and let $0\le p\le q\le n$
be such that
\[
 g_r=(a,\gamma_r)\ (r\le p),\qquad
 g_r=(\gamma'_r,\gamma_r)\ (p<r\le q),\qquad
 g_r=(\gamma'_r,a)\ (r>q),
\]
with all $\gamma_r,\gamma'_r\in C$. (Only $p\le q$ is used; it holds because the
three blocks are consecutive in this order.) Reading the two coordinates of
$(a,a)=t^{\bA^2}(g_1,\dots,g_n)$ separately gives
\begin{align}
 a&=t^{\bA}\bigl(\underbrace{a,\dots,a}_{p},\,\gamma'_{p+1},\dots,\gamma'_{n}\bigr),
   \label{eq:row1}\\
 a&=t^{\bA}\bigl(\gamma_{1},\dots,\gamma_{q},\,\underbrace{a,\dots,a}_{n-q}\bigr).
   \label{eq:row2}
\end{align}

\emph{Claim 1.} If $z_1,\dots,z_p\in a+R$ and $z_{p+1},\dots,z_n\in B$, then
$t^{\bB}(z_1,\dots,z_n)\in a+R$.

Indeed, $(a,z_r)\in R$ for $r\le p$ because $z_r\in a+R$, and
$(\gamma'_r,z_r)\in R$ for $r>p$ because $\gamma'_r\in C$ and $z_r\in B$.
Applying $t$ coordinatewise (Lemma~\ref{lem:preservation}) and using
\eqref{eq:row1} for the first coordinate gives
$\bigl(a,\,t^{\bB}(\vec z)\bigr)\in R$, i.e.\ $t^{\bB}(\vec z)\in a+R$.

\emph{Claim 2.} If $z_1,\dots,z_q\in B$ and $z_{q+1},\dots,z_n\in a+R$, then
$t^{\bB}(z_1,\dots,z_n)\in a+R$.

Indeed, $(\gamma_r,z_r)\in R$ for $r\le q$ because $\gamma_r\in C$ and
$z_r\in B$, and $(a,z_r)\in R$ for $r>q$. Applying $t$ coordinatewise and
using \eqref{eq:row2} gives $\bigl(a,\,t^{\bB}(\vec z)\bigr)\in R$.

\emph{Conclusion.} Choose any $m$ with $p\le m\le q$; such an $m$ exists because
$p\le q$. Let
\[
 f:=t(\underbrace{x_1,\dots,x_1}_{m},\underbrace{x_2,\dots,x_2}_{n-m})\in\Tm(2).
\]
For $e\in a+R$ and $b\in B$: in $f^{\bB}(e,b)$ the first $m$ arguments equal
$e$ and the rest equal $b$; since $m\ge p$, the first $p$ arguments lie in
$a+R$ and all arguments lie in $B$, so Claim~1 gives $f^{\bB}(e,b)\in a+R$.
For $b\in B$ and $e\in a+R$: in $f^{\bB}(b,e)$ the first $m$ arguments equal
$b$ and the rest equal $e$; since $m\le q$, the arguments in positions
$q+1,\dots,n$ all equal $e\in a+R$, and all arguments lie in $B$, so Claim~2
gives $f^{\bB}(b,e)\in a+R$.

Hence $a+R\abin\bB$. By Lemma~\ref{lem:center-basic}(a) it is nonempty, and it
is a proper subuniverse of $B$ because $a\notin C$. This contradicts the
hypothesis on $\bB$ together with Convention~\ref{conv:subalgebra}.
\end{proof}

\begin{definition}[Central absorption]\label{def:central-absorption}
Let $C\le\bA$. We say $C$ \emph{centrally absorbs} $\bA$, written $C\aZ\bA$,
if
\begin{enumerate}\alphlabels\tightlist
\item $C\ab\bA$, and
\item for every $a\in A\setminus C$,
      $\;(a,a)\notin\Sg_{\bA^2}\bigl((\{a\}\times C)\cup(C\times\{a\})\bigr)$.
\end{enumerate}
\end{definition}

\begin{corollary}[Zhuk's center theorem]\label{cor:zhuk-center}
Let $\bA,\bB$ be finite idempotent algebras and $R\le\bA\times\bB$ subdirect
with left center $C$. If some $t\in\Tm(k)$ is Taylor on $B$ and $\bB$ has no
proper binary absorbing subalgebra, then $C\aZ\bA$.
\end{corollary}
\begin{proof}
Condition (a) of Definition~\ref{def:central-absorption} is
Theorem~\ref{thm:center-absorbs}. For (b), note
$(\{a\}\times C)\cup(C\times\{a\})\subseteq
 (\{a\}\times C)\cup(C\times C)\cup(C\times\{a\})$, so the generated
subuniverses are nested by Lemma~\ref{lem:generation}(c), and
Theorem~\ref{thm:center-central} applies.
\end{proof}

\section{The essential doubling trick}\label{sec:doubling}

\begin{lemma}[Zhuk--Kozik essential doubling]\label{lem:doubling}
Let $n\ge1$ and let $\bA_0,\dots,\bA_{n+1}$ be algebras with $\bA_{n+1}$ finite
and idempotent. Let $C\le\bA_0$, $B_i\le\bA_i$ for $i\in[n]$, and
$C'\le\bA_{n+1}$ with $C'\aZ\bA_{n+1}$. If there exists a
$(C,B_1,\dots,B_n,C')$-essential relation
$R\le\bA_0\times\cdots\times\bA_{n+1}$, then there exists a
$(C,B_1,\dots,B_n,B_n,\dots,B_1,C)$-essential relation
\[
 R'\;\le\;\bA_0\times\cdots\times\bA_n\times\bA_n\times\cdots\times\bA_0
\]
of arity $2n+2$.
\end{lemma}
\begin{proof}
For a relation $P\le\bA_0\times\cdots\times\bA_{n+1}$ write
\[
 \beta(P):=\pi_{n+1}\bigl(P\cap(C\times B_1\times\cdots\times B_n\times A_{n+1})\bigr)\;\le\;\bA_{n+1},
\]
a subuniverse by Lemma~\ref{lem:pp}(b),(c).

\emph{Step 0: choose $R$ minimizing $|\beta(R)|$.} The family of
$(C,B_1,\dots,B_n,C')$-essential relations is nonempty by hypothesis, and
$|\beta(P)|$ is a natural number bounded by $|A_{n+1}|$; choose $R$ in the
family with $|\beta(R)|$ least, and write $B':=\beta(R)$. This uses only the
finiteness of $A_{n+1}$.

By \eqref{eq:E1} at index $n+1$, $B'\ne\varnothing$. By \eqref{eq:E2},
$B'\cap C'=\varnothing$: an element of the intersection would be the last
coordinate of a tuple of $R$ lying in $C\times B_1\times\cdots\times B_n\times C'$.
By Lemma~\ref{lem:essential-nonempty}, $C$, the $B_i$ and $C'$ are all nonempty.

\emph{Step 1: $b,b'\in B'\implies b'\in\Sg_{\bA_{n+1}}(C'\cup\{b\})$.}
Fix $b\in B'$. For each $i\in\{0,\dots,n\}$ let
$\vec x^{\,(i)}\in R\cap(C\times B_1\times\cdots\times A_i\times\cdots\times B_n\times C')$
be the witness of \eqref{eq:E1} at index $i$; note
$x^{(i)}_{n+1}\in C'$. Let $\vec y\in R\cap(C\times B_1\times\cdots\times B_n\times A_{n+1})$
have $y_{n+1}=b$, which exists because $b\in B'$. Put
\[
 R'':=\Sg_{\bA_0\times\cdots\times\bA_{n+1}}\bigl(\{\vec x^{\,(0)},\dots,\vec x^{\,(n)},\vec y\}\bigr)\subseteq R .
\]
Then $R''$ is $(C,B_1,\dots,B_n,C')$-essential: \eqref{eq:E1} at index $i\le n$
is witnessed by $\vec x^{\,(i)}$, \eqref{eq:E1} at index $n+1$ by $\vec y$, and
\eqref{eq:E2} holds because $R''\subseteq R$. By minimality
$|\beta(R'')|\ge|B'|$, while $\beta(R'')\subseteq\beta(R)=B'$ because
$R''\subseteq R$; hence $\beta(R'')=B'$. On the other hand
Lemma~\ref{lem:hom-image}(c) applied to $\pi_{n+1}$ gives
\[
 \pi_{n+1}(R'')=\Sg_{\bA_{n+1}}\bigl(\{x^{(0)}_{n+1},\dots,x^{(n)}_{n+1},b\}\bigr)
 \subseteq\Sg_{\bA_{n+1}}\bigl(C'\cup\{b\}\bigr),
\]
using Lemma~\ref{lem:generation}(c) for the inclusion. Since
$B'=\beta(R'')\subseteq\pi_{n+1}(R'')$, every $b'\in B'$ lies in
$\Sg_{\bA_{n+1}}(C'\cup\{b\})$.

\emph{Step 2: the linking relation.} Fix $b\in B'$ as in Step 1 and set
\[
 S:=\Sg_{\bA_{n+1}^2}\bigl((\{b\}\times C')\cup(C'\times\{b\})\bigr)\;\le\;\bA_{n+1}^2 .
\]
The generating set is symmetric, so $S$ is a symmetric relation; the source
records this, but the argument below never uses it, drawing instead on the two
generating blocks separately. A formalization may therefore omit the symmetry
of $S$ entirely.

\emph{Step 3: $S\cap(B'\times B')=\varnothing$.} Suppose $(b',b'')\in S$ with
$b',b''\in B'$.

Let $P:=\pi_1\bigl(S\cap(A_{n+1}\times B')\bigr)\le\bA_{n+1}$, the set of
elements having an $S$-neighbor in $B'$, a subuniverse by
Lemma~\ref{lem:pp}(b),(c). Every $c\in C'$ lies in $P$, since $(c,b)\in S$ and
$b\in B'$; and $b'\in P$, since $(b',b'')\in S$ and $b''\in B'$. Hence
$\Sg_{\bA_{n+1}}(C'\cup\{b'\})\subseteq P$ by Lemma~\ref{lem:generation}(c). By
Step 1 applied to the pair $b',b\in B'$ we have
$b\in\Sg_{\bA_{n+1}}(C'\cup\{b'\})$, so $b\in P$: there is $b'''\in B'$ with
$(b,b''')\in S$.

Let $Q:=\pi_2\bigl(S\cap(\{b\}\times A_{n+1})\bigr)\le\bA_{n+1}$, using
$\{b\}\le\bA_{n+1}$ from Lemma~\ref{lem:singleton}. Every $c\in C'$ lies in $Q$,
since $(b,c)\in S$; and $b'''\in Q$. Hence
$\Sg_{\bA_{n+1}}(C'\cup\{b'''\})\subseteq Q$. By Step~1 applied to the pair
$b''',b\in B'$ we have $b\in\Sg_{\bA_{n+1}}(C'\cup\{b'''\})$, so $b\in Q$, that
is, $(b,b)\in S$.

But $b\notin C'$, because $b\in B'$ and $B'\cap C'=\varnothing$. So
$(b,b)\in S$ contradicts condition (b) of
Definition~\ref{def:central-absorption} for $C'\aZ\bA_{n+1}$, whose generating
set is exactly the one defining $S$.

\emph{Step 4: the doubled relation.} Define
\[
 R':=\bigl\{(x_0,\dots,x_n,y_n,\dots,y_0)\;:\;
   \exists\,x_{n+1},y_{n+1}\in A_{n+1}\ \ \vec x\in R,\ \vec y\in R,\ (x_{n+1},y_{n+1})\in S\bigr\},
\]
where $\vec x=(x_0,\dots,x_{n+1})$ and $\vec y=(y_0,\dots,y_{n+1})$. It is a
subuniverse of $\bA_0\times\cdots\times\bA_n\times\bA_n\times\cdots\times\bA_0$:
form the relation
\[
 T:=\bigl\{(\vec x,\vec y)\in
 \textstyle\prod_{i=0}^{n+1}A_i\times\prod_{i=0}^{n+1}A_i \;:\;
   \vec x\in R,\ \vec y\in R,\ (x_{n+1},y_{n+1})\in S\bigr\},
\]
which is an intersection of three cylinders --- the preimages of $R$, of $R$ and
of $S$ under the projections onto the $\vec x$ coordinates, the $\vec y$
coordinates, and the pair $(x_{n+1},y_{n+1})$ respectively --- each a
subuniverse of the $(2n+4)$-fold product by Lemma~\ref{lem:pp}(d), so
$T$ is a subuniverse by Lemma~\ref{lem:pp}(a). Now project away the two
coordinates carrying $x_{n+1},y_{n+1}$ and reorder the surviving $y$
coordinates into the order $y_n,\dots,y_0$, using Lemma~\ref{lem:pp}(c),(e).

Label the coordinates of $R'$ as $0,\dots,n$ on the $x$-side and
$n^{\dagger},\dots,0^{\dagger}$ on the $y$-side, with designated subuniverses
$C,B_1,\dots,B_n$ and $B_n,\dots,B_1,C$ respectively.

\eqref{eq:E1} \emph{at an $x$-side coordinate $i\in\{0,\dots,n\}$}: take
$\vec x=\vec x^{\,(i)}$, so $x_{n+1}\in C'$, and take
$\vec y\in R\cap(C\times B_1\times\cdots\times B_n\times\{b\})$, which exists
because $b\in B'$. Then $(x_{n+1},b)\in C'\times\{b\}\subseteq S$, so the
resulting tuple lies in $R'$; its $x$-side has the free coordinate at $i$ and
designated entries elsewhere, and its $y$-side is entirely designated.

\eqref{eq:E1} \emph{at a $y$-side coordinate}: symmetric, exchanging the roles
of $\vec x$ and $\vec y$ and using $\{b\}\times C'\subseteq S$.

\eqref{eq:E2}: suppose $(x_0,\dots,x_n,y_n,\dots,y_0)\in R'$ with
$x_0,y_0\in C$ and $x_i,y_i\in B_i$ for $i\in[n]$, witnessed by
$x_{n+1},y_{n+1}$. Then $x_{n+1}\in\beta(R)=B'$ and $y_{n+1}\in B'$ by the
definition of $\beta$, while $(x_{n+1},y_{n+1})\in S$; this contradicts Step~3.
\end{proof}

\section{Central absorption is ternary}\label{sec:ternary}

\begin{corollary}[Central absorption is witnessed by a ternary term]\label{cor:central-ternary}
Let $\bA$ be a finite idempotent algebra and $C\le\bA$ with $C\aZ\bA$. Then
some $s\in\Tm(3)$ witnesses $C\ab\bA$.
\end{corollary}
\begin{proof}
Suppose not. By Theorem~\ref{thm:relational-absorption} there exists a
$C$-essential relation of arity $3$. Since $C\ab\bA$ holds by
Definition~\ref{def:central-absorption}(a),
Corollary~\ref{cor:unbounded-no-absorption} provides a largest arity $M$ of a
$C$-essential relation, and $M\ge3$.

Let $R\le\bA^{M}$ be $C$-essential of arity $M$ and write $M=n+2$ with
$n=M-2\ge1$. Then $R$ is $(C,B_1,\dots,B_n,C')$-essential for
$\bA_0=\cdots=\bA_{n+1}=\bA$, $B_1=\cdots=B_n=C$ and $C'=C$, and $C'\aZ\bA$ by
hypothesis. Lemma~\ref{lem:doubling} produces a
$(C,C,\dots,C)$-essential relation of arity $2n+2=2M-2$. Since $M\ge3$ we have
$2M-2\ge M+1>M$, contradicting the maximality of $M$.
\end{proof}

\begin{remark}[The arities produced]\label{rem:arities}
Iterating Lemma~\ref{lem:doubling} from arity $3$ produces $C$-essential
relations of arities $3,4,6,10,18,\dots$, that is $2+2^{k}$, since
$2(2+2^{k})-2=2+2^{k+1}$. Corollary~\ref{cor:central-ternary} is arranged so
that only the strict increase $2M-2>M$ for $M\ge3$ is needed, which removes the
induction.
\end{remark}

\begin{proof}[Proof of Theorem~\ref{thm:main}]
By Corollary~\ref{cor:zhuk-center}, $C\aZ\bA$. By
Corollary~\ref{cor:central-ternary} there is $s\in\Tm(3)$ witnessing
$C\ab\bA$, which by Definition~\ref{def:absorption} is exactly the displayed
containment.
\end{proof}

\begin{remark}[What is \emph{not} proved here]\label{rem:not-proved}
Section~3.10 of~\cite{BradyNotes} continues past
Corollary~\ref{cor:central-ternary}: central absorption is shown to be
compatible with quotients, transitively closed and compatible with products,
hence a ``good absorption concept''; a single ternary term is shown to witness
$C\aZ\bB$ simultaneously for all $C\aZ\bB$ with $\bB\in HSP(\bA)$; and a
criterion is given for central absorption to imply binary absorption. Those
results, and the Absorption Theorem of~\cite{BartoKozikCyclic} that they feed,
are outside the scope of this document.
\end{remark}

\appendix

\section{Statement-level citation index}\label{app:dependency-graph}

This appendix records, for every labeled numbered statement, the labeled
statements cited in its statement and proof. The table is generated
mechanically from the citation commands in the source, so it is a syntactic
citation list rather than a semantic closure: a fact used without a citation
command does not appear here. Facts imported from outside the manuscript are
collected instead in the intended imported-background package of
\hyperref[app:background]{Appendix~C}. One row cites forward: the main theorem,
Theorem~\ref{thm:main}, which is stated before its deferred proof.

\renewcommand{\arraystretch}{1.12}
\begin{longtable}{@{}>{\raggedright\arraybackslash}p{0.36\textwidth}>{\raggedright\arraybackslash}p{0.58\textwidth}@{}}
\hline
Statement & Direct internal prerequisites \\
\hline
\endfirsthead
\hline
Statement & Direct internal prerequisites \\
\hline
\endhead
\hline
\multicolumn{2}{r}{\small Continued on next page} \\
\endfoot
\hline
\endlastfoot
Theorem~\ref{thm:main}\newline\emph{Main theorem} & Definition~\ref{def:taylor}, Definition~\ref{def:absorption}, Corollary~\ref{cor:zhuk-center} and Corollary~\ref{cor:central-ternary} \\
Definition~\ref{def:algebra}\newline\emph{Signature and algebra} & Definitions and imported background (\hyperref[app:background]{Appendix~C}) only. \\
Definition~\ref{def:idempotent}\newline\emph{Idempotence} & Definitions and imported background (\hyperref[app:background]{Appendix~C}) only. \\
Definition~\ref{def:subuniverse}\newline\emph{Subuniverse} & Definitions and imported background (\hyperref[app:background]{Appendix~C}) only. \\
Convention~\ref{conv:subalgebra}\newline\emph{Subalgebra means nonempty} & Definitions and imported background (\hyperref[app:background]{Appendix~C}) only. \\
Lemma~\ref{lem:generation}\newline\emph{Intersections and generation} & Definitions and imported background (\hyperref[app:background]{Appendix~C}) only. \\
Lemma~\ref{lem:singleton}\newline\emph{Singletons} & Definition~\ref{def:idempotent} \\
Definition~\ref{def:product}\newline\emph{Products} & Definitions and imported background (\hyperref[app:background]{Appendix~C}) only. \\
Definition~\ref{def:term}\newline\emph{Terms} & Definitions and imported background (\hyperref[app:background]{Appendix~C}) only. \\
Lemma~\ref{lem:term-idempotent}\newline\emph{Term operations are idempotent} & Definition~\ref{def:idempotent} \\
Lemma~\ref{lem:preservation}\newline\emph{Preservation} & Definition~\ref{def:subuniverse} \\
Lemma~\ref{lem:sg-terms}\newline\emph{Generation by a fixed list} & Lemma~\ref{lem:preservation} and Lemma~\ref{lem:generation} \\
Remark~\ref{rem:fixed-list}\newline\emph{Why a fixed list} & Theorem~\ref{thm:center-central} \\
Definition~\ref{def:hom}\newline\emph{Homomorphism} & Definitions and imported background (\hyperref[app:background]{Appendix~C}) only. \\
Lemma~\ref{lem:hom-image}\newline\emph{Images and generation} & Lemma~\ref{lem:sg-terms} \\
Lemma~\ref{lem:pp}\newline\emph{Relational constructions} & Lemma~\ref{lem:generation} and Lemma~\ref{lem:hom-image} \\
Remark~\ref{rem:pp-usage}\newline\emph{The only relational constructions used} & Lemma~\ref{lem:pp}, Definition~\ref{def:neighborhood} and Lemma~\ref{lem:doubling} \\
Definition~\ref{def:absorption}\newline\emph{Absorption} & Definitions and imported background (\hyperref[app:background]{Appendix~C}) only. \\
Remark~\ref{rem:absorption-quantifier}\newline\emph{The quantifier form matters} & Definition~\ref{def:absorption} and Theorem~\ref{thm:relational-absorption} \\
Remark~\ref{rem:absorption-degenerate}\newline\emph{Empty and full} & Convention~\ref{conv:subalgebra} \\
Definition~\ref{def:taylor}\newline\emph{Taylor identities} & Definitions and imported background (\hyperref[app:background]{Appendix~C}) only. \\
Remark~\ref{rem:taylor-usual}\newline\emph{Relation to the usual definition} & Definition~\ref{def:taylor} and Theorem~\ref{thm:center-step} \\
Lemma~\ref{lem:taylor-binary}\newline\emph{One-sided closure gives binary absorption} & Definition~\ref{def:absorption} \\
Definition~\ref{def:star}\newline\emph{Star powers} & Definitions and imported background (\hyperref[app:background]{Appendix~C}) only. \\
Definition~\ref{def:essential}\newline\emph{Essential relation} & Definitions and imported background (\hyperref[app:background]{Appendix~C}) only. \\
Lemma~\ref{lem:essential-nonempty}\newline\emph{Essentiality forces nonemptiness} & Definitions and imported background (\hyperref[app:background]{Appendix~C}) only. \\
Proposition~\ref{prop:essential-down}\newline\emph{Passing to smaller arity} & Lemma~\ref{lem:pp} \\
Proposition~\ref{prop:absorption-blocks}\newline\emph{Absorption forbids essential relations} & Lemma~\ref{lem:preservation} and Definition~\ref{def:absorption} \\
Corollary~\ref{cor:absorption-bound}\newline\emph{Absorption bounds essential arity} & Proposition~\ref{prop:essential-down} and Proposition~\ref{prop:absorption-blocks} \\
Lemma~\ref{lem:essential-groups}\newline\emph{Regrouping} & Lemma~\ref{lem:pp} \\
Lemma~\ref{lem:clo-subuniverse}\newline\emph{The free algebra on $m$ generators} & Definitions and imported background (\hyperref[app:background]{Appendix~C}) only. \\
Theorem~\ref{thm:relational-absorption}\newline\emph{Relational description of absorption} & Proposition~\ref{prop:absorption-blocks}, Definition~\ref{def:absorption}, Lemma~\ref{lem:clo-subuniverse}, Lemma~\ref{lem:pp}, Lemma~\ref{lem:essential-groups} and Lemma~\ref{lem:preservation} \\
Corollary~\ref{cor:unbounded-no-absorption}\newline\emph{Unbounded arity defeats absorption} & Corollary~\ref{cor:absorption-bound} \\
Definition~\ref{def:subdirect}\newline\emph{Subdirect} & Definitions and imported background (\hyperref[app:background]{Appendix~C}) only. \\
Definition~\ref{def:neighborhood}\newline\emph{Neighborhood and left center} & Definitions and imported background (\hyperref[app:background]{Appendix~C}) only. \\
Lemma~\ref{lem:center-basic}\newline\emph{Basic properties} & Lemma~\ref{lem:singleton}, Lemma~\ref{lem:pp} and Lemma~\ref{lem:preservation} \\
Remark~\ref{rem:degenerate-center}\newline\emph{Degenerate centers} & Remark~\ref{rem:absorption-degenerate} \\
Theorem~\ref{thm:center-step}\newline\emph{One step towards the center} & Lemma~\ref{lem:center-basic}, Lemma~\ref{lem:term-idempotent}, Lemma~\ref{lem:taylor-binary}, Definition~\ref{def:taylor} and Convention~\ref{conv:subalgebra} \\
Theorem~\ref{thm:center-absorbs}\newline\emph{Zhuk: the left center absorbs} & Definition~\ref{def:star}, Lemma~\ref{lem:center-basic}, Lemma~\ref{lem:preservation}, Theorem~\ref{thm:center-step} and Definition~\ref{def:absorption} \\
Remark~\ref{rem:star-degenerate}\newline\emph{The degenerate arity} & Definitions and imported background (\hyperref[app:background]{Appendix~C}) only. \\
Theorem~\ref{thm:center-central}\newline\emph{Zhuk: centrality of the left center} & Lemma~\ref{lem:sg-terms}, Lemma~\ref{lem:preservation}, Lemma~\ref{lem:center-basic} and Convention~\ref{conv:subalgebra} \\
Definition~\ref{def:central-absorption}\newline\emph{Central absorption} & Definitions and imported background (\hyperref[app:background]{Appendix~C}) only. \\
Corollary~\ref{cor:zhuk-center}\newline\emph{Zhuk's center theorem} & Definition~\ref{def:central-absorption}, Theorem~\ref{thm:center-absorbs}, Lemma~\ref{lem:generation} and Theorem~\ref{thm:center-central} \\
Lemma~\ref{lem:doubling}\newline\emph{Zhuk--Kozik essential doubling} & Lemma~\ref{lem:pp}, Lemma~\ref{lem:essential-nonempty}, Lemma~\ref{lem:hom-image}, Lemma~\ref{lem:generation}, Lemma~\ref{lem:singleton} and Definition~\ref{def:central-absorption} \\
Corollary~\ref{cor:central-ternary}\newline\emph{Central absorption is witnessed by a ternary term} & Theorem~\ref{thm:relational-absorption}, Definition~\ref{def:central-absorption}, Corollary~\ref{cor:unbounded-no-absorption} and Lemma~\ref{lem:doubling} \\
Remark~\ref{rem:arities}\newline\emph{The arities produced} & Lemma~\ref{lem:doubling} and Corollary~\ref{cor:central-ternary} \\
Remark~\ref{rem:not-proved}\newline\emph{What is \emph{not} proved here} & Corollary~\ref{cor:central-ternary} \\
\emph{Narrative outside labeled statements}\newline(prose of Parts~I--III) & Theorem~\ref{thm:relational-absorption}, Lemma~\ref{lem:taylor-binary}, Lemma~\ref{lem:essential-groups}, Theorem~\ref{thm:center-step}, Theorem~\ref{thm:center-absorbs}, Theorem~\ref{thm:center-central}, Lemma~\ref{lem:doubling}, Corollary~\ref{cor:central-ternary}, Convention~\ref{conv:subalgebra}, Definition~\ref{def:central-absorption}, Lemma~\ref{lem:preservation} and Proposition~\ref{prop:absorption-blocks} \\
\end{longtable}

\section{Suggested module order}\label{app:module-order}

The citation index above, together with the imported background of
\hyperref[app:background]{Appendix~C}, approximates but does not determine the
import obligations of each module. A coarser topological order, useful for
organizing source files, is:
\begin{enumerate}\tightlist
\item signatures, algebras, subuniverses, generation, products, homomorphisms;
\item terms, term operations, preservation, generation by a fixed list,
      variable renaming along a permutation;
\item the relational constructions of Lemma~\ref{lem:pp} and the singleton
      lemma;
\item absorption, binary absorption, Taylor identities, and
      Lemma~\ref{lem:taylor-binary};
\item essential relations, arity reduction, and the block-regrouping
      Lemma~\ref{lem:essential-groups};
\item $\Clo_m(\bA)$ as a subuniverse of $\bA^{A^m}$, and the relational
      description of absorption;
\item neighborhoods $a+R$, left centers, the enlargement step, star powers, and
      Theorem~\ref{thm:center-absorbs};
\item central absorption, Theorem~\ref{thm:center-central}, the doubling trick,
      and the ternary collapse.
\end{enumerate}

Modules 1--3 are reusable universal-algebra infrastructure with no reference to
absorption. Module 6 is the only place where an algebra of functions
($\bA^{A^m}$) is formed, and hence the only place where an exponentially large
index set appears; nothing is computed there, so the size is immaterial, but a
formalization should expect the type $\bA^{A^m}$ rather than a list-based
encoding of terms to be the convenient object.

\section{Imported background}\label{app:background}

The following background package collects, in the form used, the facts imported
but not proved in this manuscript. A formal development should obtain these
from its ambient libraries.

\begin{enumerate}
\item \textbf{Finite sets and cardinality.} For finite sets $X\subseteq Y$:
      $|X|\le|Y|$, with equality if and only if $X=Y$; and if
      $X\subsetneq Y$ then $|X|+1\le|Y|$. A nonempty set of natural numbers
      bounded above has a greatest element; a nonempty set of natural numbers
      has a least element.
\item \textbf{Arithmetic of $\min$.} For natural numbers,
      $\min(N,\min(N,x)+1)=\min(N,x+1)$.
\item \textbf{Induction.} Structural induction on an inductively defined type
      of terms; strong induction on the natural numbers.
\item \textbf{Products and functions.} Indexed products of sets, tuples and
      their coordinate projections, restriction of a function to a subset,
      composition, images and preimages, and the elementary facts that images
      preserve unions and that a projection of an intersection is contained in
      the intersection of the projections.
\item \textbf{Permutations.} A finite list may be reordered by a permutation of
      its index set; equivalently, for a finite set $G$ and a partition of $G$
      into three blocks there is an enumeration of $G$ listing the blocks
      consecutively. Only this consequence is used
      (Theorem~\ref{thm:center-central}).
\end{enumerate}

The list is intended to cover the imported inferences, without a claim of
demonstrated exhaustiveness; reviewers are invited to report any use of
background not covered by an item above.

\section{Concordance with the source}\label{app:concordance}

Section~3.10 of~\cite{BradyNotes} states its results under the numbers in the
left column. The correspondence is one-to-one except that the enlargement step,
which appears inside the proof of Theorem 3.10.5 there, is isolated here.

\begin{longtable}{@{}p{0.34\textwidth}p{0.58\textwidth}@{}}
\textbf{\cite{BradyNotes}, \S3.8 and \S3.10} & \textbf{This document}\\
\hline
Proposition 3.8.2 & Proposition~\ref{prop:essential-down}\\
Proposition 3.8.3 & Proposition~\ref{prop:absorption-blocks}\\
Corollary 3.8.4 & not used\\
Theorem 3.8.5 & Theorem~\ref{thm:relational-absorption}\\
Lemma 3.8.6 & Lemma~\ref{lem:essential-groups}\\
Definition 3.10.3 & Definition~\ref{def:neighborhood}\\
Lemma 3.10.4 & Lemma~\ref{lem:taylor-binary}\\
Theorem 3.10.5 & Theorem~\ref{thm:center-step} and Theorem~\ref{thm:center-absorbs}\\
Theorem 3.10.6 & Theorem~\ref{thm:center-central}\\
Definition 3.10.7 & Definition~\ref{def:central-absorption}\\
Corollary 3.10.8 & Corollary~\ref{cor:zhuk-center}\\
Lemma 3.10.9 & Lemma~\ref{lem:doubling}\\
Corollary 3.10.10 & Corollary~\ref{cor:central-ternary}\\
Corollaries 3.10.11--3.10.13 & not covered; see Remark~\ref{rem:not-proved}\\
Propositions 3.10.14--3.10.19 & not covered; see Remark~\ref{rem:not-proved}\\
\end{longtable}

The numbering of \S3.8 above follows the printed notes; Proposition 3.8.2 is
the statement labelled \texttt{prop-essential-down} in the source, Proposition
3.8.3 is \texttt{prop-absorption-essential}, Theorem 3.8.5 is
\texttt{absorption-essential} and Lemma 3.8.6 is \texttt{essential-groups}.

\begin{thebibliography}{9}
\bibitem{BartoKozikCyclic}
L.~Barto and M.~Kozik,
\emph{Absorbing subalgebras, cyclic terms, and the constraint satisfaction problem},
Logical Methods in Computer Science \textbf{8} (2012), no.~1, 1--26.

\bibitem{BartoKozikDeciding}
L.~Barto and M.~Kozik,
\emph{Absorption in universal algebra and CSP},
in: The Constraint Satisfaction Problem: Complexity and Approximability,
Dagstuhl Follow-Ups \textbf{7} (2017), 45--77.

\bibitem{BradyNotes}
Z.~Brady,
\emph{Notes on CSPs and polymorphisms},
arXiv:2210.07383; current version at \url{https://notzeb.com/csp-notes.pdf}.

\bibitem{Zhuk}
D.~Zhuk,
\emph{A proof of the CSP dichotomy conjecture},
Journal of the ACM \textbf{67} (2020), no.~5, article 30.
\end{thebibliography}

\end{document}
